%%% FIELD reference-gradient-statement
Let \(P\) satisfy retained-cell overlap and the fourth-moment condition, and let \(P_\epsilon\) be any regular dominated observed-law submodel through \(P=P_0\) with score \(s\in L_2^0(P)\).  For
\[
\psi_P^0(O)=\sum_{k=(g,t)\in\mathcal K}\omega_k\Psi_{k,c_0(k),P}(O),
\]
the observable functional \(\tau_\omega(P)=\omega'\theta(P)\) is pathwise differentiable along the submodel and
\[
\left.\frac{d}{d\epsilon}\tau_\omega(P_\epsilon)\right|_{\epsilon=0}
=E_P[\psi_P^0(O)s(O)].
\]
In particular, \(\psi_P^0\) is a mean-zero ambient observed-law gradient of the identified aggregate.
%%% FIELD reference-gradient-proof
Fix \(k=(g,t)\) and write \(c=c_0(k)\), \(\mu_k=E_P[p_{g,1}(X)q_{k,c}(X)]\), so that \(\theta_k=\mu_k/\pi_{g,1}\).  All sums below contain at most \(T-1\) terms.  Retained-cell overlap makes every displayed inverse propensity bounded by \(\varepsilon^{-1}\), and the fourth-moment assumption and Jensen's inequality make all products used below integrable.  Thus differentiation under the finite sums and expectations is legitimate for a regular dominated submodel.

Use a dot for a derivative at zero.  Direct differentiation of a conditional probability and a conditional mean gives
\[
\dot p_{a,v}(x)=E_P[(\mathbf 1_{a,v}-p_{a,v}(x))s(O)\mid X=x]
\]
and
\[
\dot m_{a,v,u}(x)
=\frac{E_P[\mathbf 1_{a,v}\{\Delta Y_u-m_{a,v,u}(x)\}s(O)\mid X=x]}{p_{a,v}(x)}.
\]
Indeed, each identity follows by differentiating respectively the numerator and denominator of the relevant conditional expectation and applying the quotient rule; positivity of \(p_{a,v}\) is supplied by retained-cell overlap.  Differentiating \(\mu_k=E_P[p_{g,1}(X)q_{k,c}(X)]\) yields
\[
\dot\mu_k
=E_P[p_{g,1}(X)q_{k,c}(X)s(O)]
+E_P[\dot p_{g,1}(X)q_{k,c}(X)]
+E_P[p_{g,1}(X)\dot q_{k,c}(X)].
\]
The first two terms combine by iterated expectation to
\[
E_P[\mathbf 1_{g,1}q_{k,c}(X)s(O)].
\]
For example, the contribution of \(m_{g,0,u}\) to the last term is
\[
-E_P\!\left[\frac{p_{g,1}(X)}{p_{g,0}(X)}\mathbf 1_{g,0}
\{\Delta Y_u-m_{g,0,u}(X)\}s(O)\right],
\]
and the other three cell contributions follow with the signs in the definition of \(q_{k,c}\).  Summing over \(u=g,\ldots,t\), and using the residual notation in Definition \(\ref{def:pairwise-score}\), therefore gives
\[
\dot\mu_k
=E_P\!\left[\pi_{g,1}R_{k,c,P}(O)s(O)\right].
\]
Also
\[
\dot\pi_{g,1}=E_P[(\mathbf 1_{g,1}-\pi_{g,1})s(O)]
=E_P[\mathbf 1_{g,1}s(O)],
\]
where the last equality uses \(E_P[s]=0\).  Since \(\pi_{g,1}\geq\varepsilon>0\), the quotient rule now gives
\[
\dot\theta_k
=\frac{\dot\mu_k-\theta_k\dot\pi_{g,1}}{\pi_{g,1}}
=E_P\!\left[\left\{R_{k,c,P}(O)-\frac{\theta_k}{\pi_{g,1}}\mathbf 1_{g,1}\right\}s(O)\right]
=E_P[\Psi_{k,c,P}(O)s(O)].
\]
The bracket has mean zero because \(E_P[R_{k,c,P}]=\theta_k\).  Multiplication by the fixed \(\omega_k\) and summation over the finite set \(\mathcal K\) prove the assertion.
%%% FIELD complete-normal-definition
Let \(\mathscr E_0\) contain every triple \(e=(u,c,c')\) indexing an equality in Theorem \(\ref{thm:observable-equivalence}\), and associate with it the formal row
\[
\rho_e=e_{(c,1,u)}-e_{(c,0,u)}-e_{(c',1,u)}+e_{(c',0,u)}
\]
in the finite Euclidean space indexed by retained cohort--eligibility--period cells.  The deterministic set \(\mathcal E_{\mathrm{per}}\) is the lexicographically first maximal linearly independent subcollection of these rows.  If there is no nonzero restriction row, use the zero-dimensional vector and matrix convention.  For \(e=(u,c,c')\in\mathcal E_{\mathrm{per}}\), set
\[
Z^{\mathrm{per}}_{e,P}(O)=
\frac{\mathbf 1_{c,1}}{p_{c,1}(X)}\{\Delta Y_u-m_{c,1,u}(X)\}
-\frac{\mathbf 1_{c,0}}{p_{c,0}(X)}\{\Delta Y_u-m_{c,0,u}(X)\}
-\frac{\mathbf 1_{c',1}}{p_{c',1}(X)}\{\Delta Y_u-m_{c',1,u}(X)\}
+\frac{\mathbf 1_{c',0}}{p_{c',0}(X)}\{\Delta Y_u-m_{c',0,u}(X)\},
\]
and stack these coordinates as \(Z_P^{\mathrm{per}}\).  Define
\[
\psi_P^0=\sum_{k\in\mathcal K}\omega_k\Psi_{k,c_0(k),P},\qquad
K_P(X)=E_P[Z_P^{\mathrm{per}}Z_P^{\mathrm{per}\prime}\mid X],\qquad
r_P(X)=E_P[Z_P^{\mathrm{per}}\psi_P^0\mid X].
\]
The complete-normal projection, canonical-gradient candidate, and aggregate-score projection are
\[
p_{\mathrm{per},P}=r_P(X)'K_P(X)^\dagger Z_P^{\mathrm{per}},\qquad
\phi_{\mathrm{can},P}=\psi_P^0-p_{\mathrm{per},P},\qquad
\phi_{\mathrm{agg},P}=b^\star(X)'\Psi_P.
\]
Finally set \(V_{\mathrm{can}}=E_P[\phi_{\mathrm{can},P}^2]\), \(V_{\mathrm{agg}}=E_P[\phi_{\mathrm{agg},P}^2]\), and \(\Delta_{\mathrm{normal}}=V_{\mathrm{agg}}-V_{\mathrm{can}}\).
%%% FIELD complete-normal-theorem-statement
Let \(P\) be an interior observed law satisfying the observable period-specific clean-comparison equalities and admitting a full-data extension in \(\mathcal M\).  Under retained-cell overlap and the stated fourth-moment condition, the full tangent orthocomplement is
\[
\mathcal T_P^\perp
=\overline{\{a(X)'Z_P^{\mathrm{per}}(O):a\text{ measurable and }
E_P[(a(X)'Z_P^{\mathrm{per}}(O))^2]<\infty\}}^{L_2(P)}.
\]
With \(K_P\), \(r_P\), and \(\psi_P^0\) as in Definition \(\ref{def:complete-normal-projection}\),
\[
\phi_{\mathrm{can},P}(O)
=\psi_P^0(O)-r_P(X)'K_P(X)^\dagger Z_P^{\mathrm{per}}(O)
\]
is the semiparametric canonical gradient of \(\tau_\omega\), including when \(K_P(X)\) is singular or has redundant coordinates.  Its efficiency bound is
\[
V_{\mathrm{can}}
=E_P[(\psi_P^0)^2]-E_P[r_P(X)'K_P(X)^\dagger r_P(X)].
\]

The quantity \(\phi_{\mathrm{agg},P}=b^\star(X)'\Psi_P\) is instead the projection of \(\psi_P^0\) off the aggregate comparison closure \(\mathcal C_P\).  It obeys
\[
V_{\mathrm{agg}}
=\operatorname{Var}_P\{\omega'h(X)\}
+E_P[b^\star(X)'\Gamma_P(X)b^\star(X)],
\]
and every measurable feasible loading \(b\), with \(A'b(X)=\omega\) almost surely and \(b(X)'\Psi_P\in L_2(P)\), satisfies
\[
\operatorname{Var}_P\{b(X)'\Psi_P\}-V_{\mathrm{agg}}
=E_P[(b(X)-b^\star(X))'\Gamma_P(X)(b(X)-b^\star(X))]\geq0.
\]
Writing \(p_{\mathrm{agg},P}=\psi_P^0-\phi_{\mathrm{agg},P}\), the observable loss from using only aggregate normals is
\[
V_{\mathrm{agg}}-V_{\mathrm{can}}
=E_P[(p_{\mathrm{per},P}-p_{\mathrm{agg},P})^2]
=E_P[(\phi_{\mathrm{agg},P}-\phi_{\mathrm{can},P})^2]\geq0.
\]
Moreover, the following are equivalent:
\[
V_{\mathrm{agg}}=V_{\mathrm{can}},\qquad
\phi_{\mathrm{agg},P}\in\mathcal T_P,
\qquad
E_P[Z_P^{\mathrm{per}}\phi_{\mathrm{agg},P}\mid X]=0
\quad P_X\text{-almost surely}.
\]
Thus equality of the two closures is sufficient but not necessary at a particular target direction; failure of the displayed conditional projection equation gives a strictly smaller full semiparametric bound.
%%% FIELD complete-normal-theorem-proof
For each retained cell and period define the standardized residual
\[
U_{a,s,u}=\frac{\mathbf 1_{a,s}}{p_{a,s}(X)}\{\Delta Y_u-m_{a,s,u}(X)\}.
\]
Retained-cell overlap and the fourth-moment assumption imply \(U_{a,s,u}\in L_2(P)\), and the definition of \(m_{a,s,u}\) gives \(E_P[U_{a,s,u}\mid X]=0\).  Every period-specific normal in Theorem \(\ref{thm:observable-equivalence}\) is the linear image of its formal row \(\rho_e\), namely
\[
Z_e=U_{c,1,u}-U_{c,0,u}-U_{c',1,u}+U_{c',0,u}.
\]
The rows selected by \(\mathcal E_{\mathrm{per}}\) span all rows \(\rho_e\).  Hence every \(Z_e\) is a deterministic linear combination of coordinates of \(Z_P^{\mathrm{per}}\), while every coordinate of \(Z_P^{\mathrm{per}}\) is itself one of the normals in that theorem.  Multiplication by arbitrary measurable functions of \(X\), followed by \(L_2(P)\) closure, preserves both inclusions.  The tangent-exhaustion conclusion of Theorem \(\ref{thm:observable-equivalence}\) therefore gives exactly
\[
\mathcal T_P^\perp
=\overline{\{a(X)'Z_P^{\mathrm{per}}:a(X)'Z_P^{\mathrm{per}}\in L_2(P)\}}^{L_2(P)}.
\tag{1}
\]
This is a consequence of the observable restrictions; no additional spanning premise has been imposed.

The reference-score-gradient lemma shows that \(\psi_P^0\in L_2^0(P)\) represents the pathwise derivative of \(\tau_\omega\) on every score in \(\mathcal T_P\).  We next calculate its orthogonal projection onto the right side of (1).  Conditional on \(X=x\), \(K_P(x)\) is symmetric and positive semidefinite.  If \(v\in\ker K_P(x)\), then
\[
0=v'K_P(x)v=E_P[(v'Z_P^{\mathrm{per}})^2\mid X=x],
\]
so \(v'Z_P^{\mathrm{per}}=0\) conditionally almost surely and therefore
\[
v'r_P(x)=E_P[v'Z_P^{\mathrm{per}}\psi_P^0\mid X=x]=0.
\]
Thus \(r_P(x)\in\{\ker K_P(x)\}^\perp=\operatorname{range}K_P(x)\), and consequently
\[
K_P(x)K_P(x)^\dagger r_P(x)=r_P(x).
\tag{2}
\]
For a symmetric positive semidefinite matrix,
\[
K_P(x)^\dagger=\lim_{j\to\infty}\{K_P(x)^2+j^{-1}I\}^{-1}K_P(x).
\]
Every map in the sequence is continuous in the matrix entries, so the pointwise limit is Borel measurable.  Thus the coefficient \(K_P^\dagger r_P\) is measurable even if rank is not constant.

Put \(p_{\mathrm{per},P}=r_P'K_P^\dagger Z_P^{\mathrm{per}}\).  By (2), symmetry of \(K_P^\dagger\), and the Moore--Penrose identities,
\[
E_P[p_{\mathrm{per},P}^2\mid X]
=r_P'K_P^\dagger K_PK_P^\dagger r_P
=r_P'K_P^\dagger r_P
=E_P[\psi_P^0p_{\mathrm{per},P}\mid X].
\tag{3}
\]
It follows that
\[
E_P[(\psi_P^0)^2\mid X]
=E_P[(\psi_P^0-p_{\mathrm{per},P})^2\mid X]
+E_P[p_{\mathrm{per},P}^2\mid X].
\tag{4}
\]
In particular, \(p_{\mathrm{per},P}\in L_2(P)\), so it belongs to the generating set in (1).  For any measurable \(a\) for which \(a'Z_P^{\mathrm{per}}\in L_2(P)\), (2) gives
\[
E_P[(\psi_P^0-p_{\mathrm{per},P})a(X)'Z_P^{\mathrm{per}}]
=E_P[a(X)'\{r_P(X)-K_P(X)K_P(X)^\dagger r_P(X)\}]=0.
\tag{5}
\]
Equations (1), (4), and (5) prove that \(p_{\mathrm{per},P}=\Pi_{\mathcal T_P^\perp}\psi_P^0\) and
\[
\phi_{\mathrm{can},P}=\psi_P^0-p_{\mathrm{per},P}=\Pi_{\mathcal T_P}\psi_P^0.
\]
For every tangent score \(s\in\mathcal T_P\), (5) implies
\[
E_P[\phi_{\mathrm{can},P}s]=E_P[\psi_P^0s]=\dot\tau_{\omega,P}[s].
\]
Thus \(\phi_{\mathrm{can},P}\) is a gradient lying in the tangent space and hence is the canonical gradient.  Taking expectations in (3)--(4) proves the displayed formula for \(V_{\mathrm{can}}\).  Notice that singularity creates no unidentified coefficient: (2) makes the fitted random variable unique, while two coefficient solutions differ only by an element of \(\ker K_P(X)\), whose product with \(Z_P^{\mathrm{per}}\) is zero.

It remains to separate the aggregate score projection from the full projection.  On the observable equality model, the score definition and zero conditional means of the residuals give, for every \((k,c)\in\mathcal I\),
\[
E_P[\Psi_{k,c,P}\mid X]
=\frac{p_{g,1}(X)}{\pi_{g,1}}\{q_{k,c}(X)-\theta_k\}=h_k(X).
\]
In stacked notation,
\[
E_P[\Psi_P\mid X]=Ah(X).
\tag{6}
\]
Consequently every feasible \(b\) has conditional mean \(\omega'h(X)\), and the law of total variance yields
\[
\operatorname{Var}_P\{b(X)'\Psi_P\}
=\operatorname{Var}_P\{\omega'h(X)\}+E_P[b(X)'\Gamma_P(X)b(X)].
\tag{7}
\]
The singular-operator lemma gives \(A'b^\star=\omega\) and \(N'\Gamma b^\star=0\).  If \(d=b-b^\star\), then \(A'd=0\), so \(d=Nz\) for a measurable \(z\).  Hence
\[
d'\Gamma b^\star=z'N'\Gamma b^\star=0,
\]
and expansion of \(b=b^\star+d\) in (7) proves both the formula for \(V_{\mathrm{agg}}\) and the stated nonnegative excess-variance identity.

Let \(b^{\mathrm{ref}}\) put weight \(\omega_k\) on coordinate \((k,c_0(k))\) and zero on the other coordinates.  Then \(A'b^{\mathrm{ref}}=\omega\) and \(\psi_P^0=b^{\mathrm{ref}\prime}\Psi_P\).  Because \(b^{\mathrm{ref}}-b^\star\in\ker A'=\operatorname{col}N\),
\[
p_{\mathrm{agg},P}:=\psi_P^0-\phi_{\mathrm{agg},P}
=(b^{\mathrm{ref}}-b^\star(X))'\Psi_P
=(b^{\mathrm{ref}}-b^\star(X))'R_P\in\mathcal C_P,
\tag{8}
\]
where the last equality uses that the denominator correction is constant inside a target block and is annihilated by \(N'\).  Moreover, (6), \(N'A=0\), and the normal equation give
\[
E_P[N'R_P\phi_{\mathrm{agg},P}\mid X]
=E_P[N'\Psi_P\Psi_P'b^\star\mid X]
=N'\{\Gamma_P+Ahh'A'\}b^\star
=N'\Gamma_Pb^\star=0.
\tag{9}
\]
Equations (8)--(9) prove that \(p_{\mathrm{agg},P}=\Pi_{\mathcal C_P}\psi_P^0\) and \(\phi_{\mathrm{agg},P}=\psi_P^0-p_{\mathrm{agg},P}\).

By Theorem \(\ref{thm:observable-equivalence}\), \(\mathcal C_P\subseteq\mathcal T_P^\perp\).  Therefore
\[
w_P:=p_{\mathrm{per},P}-p_{\mathrm{agg},P}
\in\mathcal T_P^\perp\cap\mathcal C_P^\perp,
\qquad
\phi_{\mathrm{agg},P}=\phi_{\mathrm{can},P}+w_P.
\tag{10}
\]
Since \(\phi_{\mathrm{can},P}\in\mathcal T_P\), the two terms on the right of (10) are orthogonal.  Pythagoras gives
\[
V_{\mathrm{agg}}-V_{\mathrm{can}}=E_P[w_P^2]
=E_P[(p_{\mathrm{per},P}-p_{\mathrm{agg},P})^2]
=E_P[(\phi_{\mathrm{agg},P}-\phi_{\mathrm{can},P})^2].
\]
Equality holds exactly when \(w_P=0\), equivalently when \(\phi_{\mathrm{agg},P}\in\mathcal T_P\).  Finally, (1) implies
\[
\phi_{\mathrm{agg},P}\in\mathcal T_P
\quad\Longleftrightarrow\quad
E_P[Z_P^{\mathrm{per}}\phi_{\mathrm{agg},P}\mid X]=0
\quad P_X\text{-almost surely},
\]
because vanishing conditional covariance is equivalent to orthogonality against every square-integrable measurable multiple.  This proves all assertions.
%%% FIELD complete-normal-rates-statement
Suppose \(X\) has fixed finite support and every support mass is bounded below.  Use training-cell averages for all conditional moments in Definition \(\ref{def:crossfit-estimator}\), and suppose the propensity and regression first stages are parametric cell averages.  Under retained-cell overlap, the conditional fourth-moment envelope, fixed-fold cross-fitting, \(\tau_n\downarrow0\), and \(n^{-1/2}/\tau_n\to0\), uniformly over the fixed number of folds,
\[
\|\widehat K^{-\ell}-K_P\|_{P_X,4,\mathrm{op}}
+\|\widehat r^{-\ell}-r_P\|_{P_X,4}=O_P(n^{-1/2}).
\]
For each support point, thresholding recovers \(\operatorname{rank}K_P(x)\) with probability tending to one.  On the recovered ranges,
\[
\|\{\widehat K^{-\ell}_\tau\}^{\dagger}\widehat r^{-\ell}
-K_P^\dagger r_P\|_{P_X,4}=O_P(n^{-1/2}),
\]
while the unconditional conservative threshold bound is \(O_P(n^{-1/2}/\tau_n)\).  Consequently, with \(\tau_n=n^{-a}\), \(0<a<1/4\), the complete-normal learned-projection rate is \(O_P(n^{-1/2+a})\), its square is \(o(n^{-1/2})\), and its product with root-\(n\) propensity and regression rates is \(o(n^{-1/2})\).
%%% FIELD complete-normal-rates-proof
There are finitely many covariate support points, retained cells, periods, normal coordinates, target coordinates, and folds.  The lower support-mass bound and retained-cell overlap imply that every training count used by a conditional cell average is of order \(n\) in probability.  Each coordinate of \(Z_P^{\mathrm{per}}\) is a finite sum of outcome residuals multiplied by inverse propensities bounded by \(\varepsilon^{-1}\).  Each coordinate of \(\psi_P^0\) is likewise a finite sum of such residuals and bounded propensity ratios, plus a target-cell term.  The conditional fourth-moment envelope, fixed \(T\), and Jensen's inequality therefore yield, uniformly over support points,
\[
E_P[\|Z_P^{\mathrm{per}}\|^4+|\psi_P^0|^4\mid X]\le C_P<\infty.
\tag{11}
\]
The constant may depend on the fixed law, which is appropriate for the asserted pointwise rate.

By (11), every coordinate of \(Z_P^{\mathrm{per}}Z_P^{\mathrm{per}\prime}\) and \(Z_P^{\mathrm{per}}\psi_P^0\) has finite conditional variance.  Chebyshev's inequality gives an \(O_P(n^{-1/2})\) error for each oracle training-cell mean.  A union bound over the fixed finite collection gives the same rate jointly.  On the overlap set, the maps from \(p\), \(m\), \(\pi\), and \(\theta\) to \(Z^{\mathrm{per}}\) and \(\psi^0\) are finite sums and products of functions whose denominators are at least \(\varepsilon\).  The mean-value theorem and Hölder's inequality therefore show that replacing the oracle nuisances by root-\(n\) cell averages adds \(O_P(n^{-1/2})\) in \(L_4(P_X)\).  The preliminary ratios \(\widetilde\theta_k\) are also root-\(n\), because their numerators are finite cell means and their denominators satisfy \(\pi_{g,1}\ge\varepsilon\).  Hence
\[
\max_\ell\|\widehat K^{-\ell}-K_P\|_{P_X,4,\mathrm{op}}
+\max_\ell\|\widehat r^{-\ell}-r_P\|_{P_X,4}=O_P(n^{-1/2}).
\tag{12}
\]

At a fixed support point \(x\), let \(\kappa_P(x)\) be the smallest positive eigenvalue of \(K_P(x)\), with \(\kappa_P(x)=+\infty\) if its rank is zero.  Because the support is finite,
\[
\underline\kappa_P:=\min_{x:\,P_X(x)>0}\kappa_P(x)>0.
\tag{13}
\]
This is a consequence at the fixed law, not a new uniform model restriction.  Weyl's inequality and (12) show that, once \(\tau_n<\underline\kappa_P/2\), a rank error requires
\[
\|\widehat K^{-\ell}(x)-K_P(x)\|_{\mathrm{op}}>\tau_n.
\]
The fourth-power Markov inequality, (12), and \(n^{-1/2}/\tau_n\to0\) make the probability of this event tend to zero uniformly over the finite support and folds.  This proves rank recovery.

On the rank-recovery event, all retained eigenvalues exceed \(\underline\kappa_P/2\).  The finite-dimensional resolvent identity on the common spectral range and (12) then give
\[
\|\{\widehat K^{-\ell}_\tau\}^\dagger\widehat r^{-\ell}-K_P^\dagger r_P\|
\le C_P\{\|\widehat K^{-\ell}-K_P\|_{\mathrm{op}}
+\|\widehat r^{-\ell}-r_P\|\}=O_P(n^{-1/2}).
\tag{14}
\]
Without conditioning on correct selection, the retained inverse has norm at most \(\tau_n^{-1}\).  Splitting into the correct-selection event and its complement, and using the same deterministic perturbation decomposition as in the aggregate loading lemma, gives the conservative bound \(O_P(n^{-1/2}/\tau_n)\).  If \(\tau_n=n^{-a}\), this is \(O_P(n^{-1/2+a})\), and
\[
n^{-1+2a}=o(n^{-1/2})\quad\text{for }a<1/4,
\qquad
n^{-1/2+a}n^{-1/2}=o(n^{-1/2})\quad\text{for }a<1/2.
\]
The declared range \(0<a<1/4\) satisfies both requirements, proving the final assertions.
%%% FIELD generic-tangent-projection-statement
For any \(v\in L_2(P)\), put \(v^0=v-E_P[v]\), \(r_v(X)=E_P[Z_P^{\mathrm{per}}v^0\mid X]\), and
\[
t_v=v^0-r_v(X)'K_P(X)^\dagger Z_P^{\mathrm{per}}.
\]
At an interior law covered by Theorem \(\ref{thm:complete-normal-canonical-gradient}\), \(t_v=\Pi_{\mathcal T_P}v^0\).  In particular, every population coordinate of the complete-normal projected dictionary in Definition \(\ref{def:gap-one-step}\) belongs to \(\mathcal T_P\).
%%% FIELD generic-tangent-projection-proof
The argument is the conditional projection calculation in the complete-normal canonical-gradient theorem with \(v^0\) in place of \(\psi_P^0\).  Explicitly, if \(a\in\ker K_P(X)\), then \(a'Z_P^{\mathrm{per}}=0\) conditionally almost surely, so \(a'r_v(X)=0\).  Hence \(r_v(X)\in\operatorname{range}K_P(X)\) and \(K_PK_P^\dagger r_v=r_v\).  Therefore, for every measurable \(c\) such that \(c(X)'Z_P^{\mathrm{per}}\in L_2(P)\),
\[
E_P\!\left[\{v^0-r_v'K_P^\dagger Z_P^{\mathrm{per}}\}
c(X)'Z_P^{\mathrm{per}}\right]
=E_P[c(X)'\{r_v-K_PK_P^\dagger r_v\}]=0.
\]
The same Pythagorean calculation as in equation \((4)\) of that theorem shows that \(r_v'K_P^\dagger Z_P^{\mathrm{per}}\in L_2(P)\).  Because measurable multiples of \(Z_P^{\mathrm{per}}\) have closure \(\mathcal T_P^\perp\), the subtracted term is \(\Pi_{\mathcal T_P^\perp}v^0\), and the residual is \(\Pi_{\mathcal T_P}v^0\), as claimed.
%%% FIELD corrected-crossfit-definition
Fix folds \((I_\ell)_{\ell=1}^L\).  On the training complement of fold \(I_\ell\), estimate every retained propensity \(p_{a,s}\), regression \(m_{a,s,u}\), target share \(\pi_{g,1}\), and a preliminary target vector \(\widetilde\theta^{-\ell}\).  On a validation observation define \(\widehat R_i^{-\ell}\) from Definition \(\ref{def:pairwise-score}\), define the denominator-corrected vector \(\widehat\Psi_i^{-\ell}\) using \(\widetilde\theta^{-\ell}\), and put
\[
\widehat\psi_i^{0,-\ell}=\sum_{k\in\mathcal K}\omega_k
\widehat\Psi_{i,k,c_0(k)}^{-\ell}.
\]
Using the same training nuisances, form every coordinate of \(\widehat Z_i^{\mathrm{per},-\ell}\) from Definition \(\ref{def:complete-normal-projection}\).  Fit on the training complement the conditional moments
\[
\widehat K^{-\ell}(x)\approx E_P[Z_P^{\mathrm{per}}Z_P^{\mathrm{per}\prime}\mid X=x],
\qquad
\widehat r^{-\ell}(x)\approx E_P[Z_P^{\mathrm{per}}\psi_P^0\mid X=x].
\]
Symmetrize \(\widehat K^{-\ell}\), retain its eigenvalues above the deterministic threshold \(\tau_n\), and denote the resulting Moore--Penrose inverse by \(\{\widehat K_\tau^{-\ell}(x)\}^{\dagger}\).  The validation-fold estimated canonical influence value is
\[
\widehat\phi_{\mathrm{can},i}^{-\ell}
=\widehat\psi_i^{0,-\ell}
-\widehat r^{-\ell}(X_i)'\{\widehat K_\tau^{-\ell}(X_i)\}^{\dagger}
\widehat Z_i^{\mathrm{per},-\ell}.
\]
The cross-fitted one-step estimator, variance estimator, and Wald interval are
\[
\widehat\tau_{\mathrm{eff}}
=\sum_{\ell=1}^L\frac{|I_\ell|}{n}
\left\{\omega'\widetilde\theta^{-\ell}
+\frac{1}{|I_\ell|}\sum_{i\in I_\ell}\widehat\phi_{\mathrm{can},i}^{-\ell}\right\},
\]
\[
\widehat V_{\mathrm{eff}}
=\frac1n\sum_{\ell=1}^L\sum_{i\in I_\ell}
\left(\widehat\phi_{\mathrm{can},i}^{-\ell}-\overline{\widehat\phi}_{\mathrm{can}}\right)^2,
\qquad
\mathrm{CI}_{1-\alpha}
=\widehat\tau_{\mathrm{eff}}\pm z_{1-\alpha/2}
\sqrt{\widehat V_{\mathrm{eff}}/n}.
\]
The aggregate loading \(\widehat b^{-\ell}\) remains available for the comparator diagnostic but is not substituted for the complete-normal correction in this estimator.
%%% FIELD corrected-gap-definition
For a declared constant comparator \(b^{\mathrm{cmp}}\), define its loss relative to the full information bound by
\[
\Delta_{\mathrm{cmp}}(P)
=\operatorname{Var}_P\{b^{\mathrm{cmp}\prime}\Psi_P(O)\}-V_{\mathrm{can}}
=E_P[(b^{\mathrm{cmp}}-b^\star(X))'\Gamma_P(X)(b^{\mathrm{cmp}}-b^\star(X))]
+\Delta_{\mathrm{normal}}.
\]
For a mean-zero observed score \(s\), obtain \(D\Delta_{\mathrm{cmp},P}[s]\) by differentiating the displayed observed-law functional: differentiate \(p\), \(m\), \(\pi\), \(\theta\), and the denominator-corrected \(\Psi\) in Definition \(\ref{def:pairwise-score}\); then differentiate \(\Gamma=\operatorname{Cov}(\Psi\mid X)\), \(K=E[Z^{\mathrm{per}}Z^{\mathrm{per}\prime}\mid X]\), \(r=E[Z^{\mathrm{per}}\psi^0\mid X]\), and each locally fixed-rank Moore--Penrose map by the product and quotient rules.  This defines an observable linear derivative on every locally fixed-rank finite-support submodel.

Let \(v_{J_n}=(v_1,\ldots,v_{J_n})'\) be the declared bounded dictionary.  On each training complement, estimate \(Z^{\mathrm{per}}\), \(K\), and \(E[Z^{\mathrm{per}}\{v_j-Ev_j\}\mid X]\), and set
\[
\widehat a_j^{-\ell}(x)=\{\widehat K_\tau^{-\ell}(x)\}^{\dagger}
\widehat E_{-\ell}[\widehat Z^{\mathrm{per},-\ell}\{v_j-\mathbb E_{-\ell}v_j\}\mid X=x],
\]
\[
\widehat t_j^{-\ell}=v_j-\mathbb E_{-\ell}v_j
-\widehat a_j^{-\ell}(X)'\widehat Z^{\mathrm{per},-\ell},
\]
followed by training-fold centering.  Evaluate the preceding derivative at \(s=\widehat t_j^{-\ell}\) to form \(\widehat g_{\Delta,J_n}^{-\ell}\), set
\[
\widehat G_{J_n}^{-\ell}=\mathbb E_{-\ell}[\widehat t_{J_n}^{-\ell}
\widehat t_{J_n}^{-\ell\prime}],
\qquad
\widehat\zeta_{\Delta,J_n}^{-\ell}
=\widehat t_{J_n}^{-\ell\prime}
\{\widehat G_{J_n}^{-\ell}\}^{\dagger}\widehat g_{\Delta,J_n}^{-\ell}.
\]
The foldwise plug-in is the empirical version of the first display, using the estimated comparator, \(\widehat\Gamma^{-\ell}\), \(\widehat b^{-\ell}\), and the complete-normal estimate of \(\Delta_{\mathrm{normal}}\).  Add the validation-fold mean of \(\widehat\zeta_{\Delta,J_n}^{-\ell}\) and aggregate folds with weights \(|I_\ell|/n\).  This construction projects the Riesz dictionary against the complete period-specific normal vector; asymptotic normality is asserted only by a separate theorem that verifies local fixed-rank differentiation and the required first-stage and Riesz-map rates.
